% GENERATED FILE. Edit canonical lessons, then run npm run generate:books.
% canonical-source-hash: fnv1a64:7fdea1c851c86a27
% canonical-lessons: TA-C297-katattu, TA-C297-tappiyotu, TA-C297-kallanottu, TA-C297-pulanayvu, TA-C297-kurrancattu

\chapter{Smugglers on the Run}
\label{ch:ta-smugglers-on-the-run}
\hlchaptermodality{297}

\begin{chapteropening}
\textbf{By the end of this chapter:} \emph{I can tell a story about smugglers and an investigation.}

Complete the last lesson of chapter 297: I can tell a story about smugglers and an investigation.
\end{chapteropening}

\section[\texorpdfstring{\ta{கடத்து} (kaṭattu)}{கடத்து (kaṭattu)}]{\ta{கடத்து} (kaṭattu) — to smuggle; to kidnap}

\label{lesson:TA-C297-katattu}



\begin{quote}
Before the new one: say the Tamil for a prisoner, then the Tamil for to confess.
\end{quote}

\subsection*{You'll want to know: \ta{கடத்து}}
\textbf{\ta{கடத்து}} — \emph{kaṭattu} — \textquotedblleft{}to smuggle; to kidnap\textquotedblright{}.

\begin{quote}
\ta{அன்று} \ta{காவலர்கள்} \ta{வராமல்} \ta{இருந்திருந்தால்}, \ta{அவர்கள்} \ta{தங்கத்தைக்} \ta{கடத்தியிருப்பார்கள்}.
\end{quote}

\begin{quote}
\emph{aṉṟu kāvalarkaḷ varāmal iruntiruntāl, avarkaḷ taṅkattaik kaṭattiyiruppārkaḷ.}
\end{quote}

\begin{quote}
If the police had not come that day, they would have smuggled the gold.
\end{quote}

Past (I …): \ta{கடத்தினேன்}, \emph{kaṭattiṉēṉ}.

What is carried off takes -\ta{ஐ}; also to kidnap: \ta{குழந்தையைக்} \ta{கடத்தினார்கள்}.

\begin{cousinweb}
\textbf{\ta{கடத்து}} (\emph{kaṭattu}) is native Dravidian: DEDR 1109 gives it as the causative of \ta{கட}, \textquotedblleft{}to cross\textquotedblright{}, so to smuggle is to make something cross.
\end{cousinweb}

\begin{grammarlens}[title={What you've built}]
One more word for telling what happened.
\end{grammarlens}

\subsection*{Your turn}
\begin{itemize}
\raggedright
  \item \emph{Say it:} \emph{kaṭattu}
  \item \emph{Say it:} \emph{kaṭattu}, once more
  \item \emph{Say it:} \emph{aṉṟu kāvalarkaḷ varāmal iruntiruntāl, avarkaḷ taṅkattaik kaṭattiyiruppārkaḷ.} — If the police had not come that day, they would have smuggled the gold.
  \item \emph{Recall:} say the Tamil for a prisoner, then the Tamil for to confess, then say \emph{kaṭattu} again
\end{itemize}

\begin{tcolorbox}[breakable,colback=teal!4,colframe=teal!35!black,title={Before you move on}]
What does \ta{கடத்து} mean? (\textquotedblleft{}To smuggle; to kidnap\textquotedblright{}.) Say it once more.
\end{tcolorbox}
\section[\texorpdfstring{\ta{தப்பியோடு} (tappiyōṭu)}{தப்பியோடு (tappiyōṭu)}]{\ta{தப்பியோடு} (tappiyōṭu) — to flee}

\label{lesson:TA-C297-tappiyotu}



\begin{quote}
Before the new one: say the Tamil for to confess, then the Tamil for to smuggle.
\end{quote}

\subsection*{You'll want to know: \ta{தப்பியோடு}}
\textbf{\ta{தப்பியோடு}} — \emph{tappiyōṭu} — \textquotedblleft{}to flee\textquotedblright{}.

\begin{quote}
\ta{திருடர்கள்} \ta{இருட்டில்} \ta{தப்பியோடினார்கள்}.
\end{quote}

\begin{quote}
\emph{tiruṭarkaḷ iruṭṭil tappiyōṭiṉārkaḷ.}
\end{quote}

\begin{quote}
The thieves fled in the dark.
\end{quote}

Past (I …): \ta{தப்பியோடினேன்}, \emph{tappiyōṭiṉēṉ}.

Built from \ta{தப்பி}, having escaped, and \ta{ஓடு}, to run.

\begin{cousinweb}
\textbf{\ta{தப்பியோடு}} (\emph{tappiyōṭu}) is a Tamil compound built from \ta{தப்பு}, \textquotedblleft{}to slip away, escape\textquotedblright{} (DEDR 3071), and \ta{ஓடு}, \textquotedblleft{}to run\textquotedblright{} (DEDR 1041).
\end{cousinweb}

\begin{grammarlens}[title={What you've built}]
One more word for telling what happened.
\end{grammarlens}

\subsection*{Your turn}
\begin{itemize}
\raggedright
  \item \emph{Say it:} \emph{tappiyōṭu}
  \item \emph{Say it:} \emph{tappiyōṭu}, once more
  \item \emph{Say it:} \emph{tiruṭarkaḷ iruṭṭil tappiyōṭiṉārkaḷ.} — The thieves fled in the dark.
  \item \emph{Recall:} say the Tamil for to confess, then the Tamil for to smuggle, then say \emph{tappiyōṭu} again
\end{itemize}

\begin{tcolorbox}[breakable,colback=teal!4,colframe=teal!35!black,title={Before you move on}]
What does \ta{தப்பியோடு} mean? (\textquotedblleft{}To flee\textquotedblright{}.) Say it once more.
\end{tcolorbox}
\section[\texorpdfstring{\ta{கள்ளநோட்டு} (kaḷḷanōṭṭu)}{கள்ளநோட்டு (kaḷḷanōṭṭu)}]{\ta{கள்ளநோட்டு} (kaḷḷanōṭṭu) — a counterfeit note}

\label{lesson:TA-C297-kallanottu}



\begin{quote}
Before the new one: say the Tamil for to smuggle, then the Tamil for to flee.
\end{quote}

\subsection*{You'll want to know: \ta{கள்ளநோட்டு}}
\textbf{\ta{கள்ளநோட்டு}} — \emph{kaḷḷanōṭṭu} — \textquotedblleft{}a counterfeit note\textquotedblright{}.

\begin{quote}
\ta{கடைக்காரர்} \ta{ஒரு} \ta{கள்ளநோட்டைக்} \ta{கண்டுபிடித்தார்}.
\end{quote}

\begin{quote}
\emph{kaṭaikkārar oru kaḷḷanōṭṭaik kaṇṭupiṭittār.}
\end{quote}

\begin{quote}
The shopkeeper spotted a counterfeit note.
\end{quote}

\ta{கள்ள}- means false or stolen before a noun.

\begin{cousinweb}
\textbf{\ta{கள்ளநோட்டு}} (\emph{kaḷḷanōṭṭu}) joins native \ta{கள்ளம்}, \textquotedblleft{}deceit, theft\textquotedblright{} (DEDR 1372), with \ta{நோட்டு}, a loan from English \emph{note}, \textquotedblleft{}a banknote\textquotedblright{}.
\end{cousinweb}

\begin{grammarlens}[title={What you've built}]
One more word for telling what happened.
\end{grammarlens}

\subsection*{Your turn}
\begin{itemize}
\raggedright
  \item \emph{Say it:} \emph{kaḷḷanōṭṭu}
  \item \emph{Say it:} \emph{kaḷḷanōṭṭu}, once more
  \item \emph{Say it:} \emph{kaṭaikkārar oru kaḷḷanōṭṭaik kaṇṭupiṭittār.} — The shopkeeper spotted a counterfeit note.
  \item \emph{Recall:} say the Tamil for to smuggle, then the Tamil for to flee, then say \emph{kaḷḷanōṭṭu} again
\end{itemize}

\begin{tcolorbox}[breakable,colback=teal!4,colframe=teal!35!black,title={Before you move on}]
What does \ta{கள்ளநோட்டு} mean? (\textquotedblleft{}A counterfeit note\textquotedblright{}.) Say it once more.
\end{tcolorbox}
\section[\texorpdfstring{\ta{புலனாய்வு} (pulaṉāyvu)}{புலனாய்வு (pulaṉāyvu)}]{\ta{புலனாய்வு} (pulaṉāyvu) — an investigation}

\label{lesson:TA-C297-pulanayvu}



\begin{quote}
Before the new one: say the Tamil for to flee, then the Tamil for a counterfeit note.
\end{quote}

\subsection*{You'll want to know: \ta{புலனாய்வு}}
\textbf{\ta{புலனாய்வு}} — \emph{pulaṉāyvu} — \textquotedblleft{}an investigation\textquotedblright{}.

\begin{quote}
\ta{இந்தக்} \ta{கொள்ளை} \ta{பற்றிப்} \ta{புலனாய்வு} \ta{தொடங்கியது}.
\end{quote}

\begin{quote}
\emph{intak koḷḷai paṟṟip pulaṉāyvu toṭaṅkiyatu.}
\end{quote}

\begin{quote}
An investigation into this robbery began.
\end{quote}

News reports say \ta{புலனாய்வு} \ta{நடக்கிறது}, an investigation is under way.

\begin{cousinweb}
\textbf{\ta{புலனாய்வு}} (\emph{pulaṉāyvu}) is a Tamil compound built from \ta{புலன்}, \textquotedblleft{}the senses, perception\textquotedblright{}, and \ta{ஆய்வு}, \textquotedblleft{}examination\textquotedblright{}, from \ta{ஆய்}, \textquotedblleft{}to search, examine\textquotedblright{} (DEDR 363).
\end{cousinweb}

\begin{grammarlens}[title={What you've built}]
One more word for telling what happened.
\end{grammarlens}

\subsection*{Your turn}
\begin{itemize}
\raggedright
  \item \emph{Say it:} \emph{pulaṉāyvu}
  \item \emph{Say it:} \emph{pulaṉāyvu}, once more
  \item \emph{Say it:} \emph{intak koḷḷai paṟṟip pulaṉāyvu toṭaṅkiyatu.} — An investigation into this robbery began.
  \item \emph{Recall:} say the Tamil for to flee, then the Tamil for a counterfeit note, then say \emph{pulaṉāyvu} again
\end{itemize}

\begin{tcolorbox}[breakable,colback=teal!4,colframe=teal!35!black,title={Before you move on}]
What does \ta{புலனாய்வு} mean? (\textquotedblleft{}An investigation\textquotedblright{}.) Say it once more.
\end{tcolorbox}
\section[\texorpdfstring{\ta{குற்றஞ்சாட்டு} (kuṟṟañcāṭṭu)}{குற்றஞ்சாட்டு (kuṟṟañcāṭṭu)}]{\ta{குற்றஞ்சாட்டு} (kuṟṟañcāṭṭu) — to accuse}

\label{lesson:TA-C297-kurrancattu}



\begin{quote}
Before the new one: say the Tamil for a counterfeit note, then the Tamil for an investigation.
\end{quote}

\subsection*{You'll want to know: \ta{குற்றஞ்சாட்டு}}
\textbf{\ta{குற்றஞ்சாட்டு}} — \emph{kuṟṟañcāṭṭu} — \textquotedblleft{}to accuse\textquotedblright{}.

\begin{quote}
\ta{அவர்} \ta{என்னைப்} \ta{பொய்யாகக்} \ta{குற்றஞ்சாட்டினார்}.
\end{quote}

\begin{quote}
\emph{avar eṉṉaip poyyākak kuṟṟañcāṭṭiṉār.}
\end{quote}

\begin{quote}
He accused me falsely.
\end{quote}

Past (I …): \ta{குற்றஞ்சாட்டினேன்}, \emph{kuṟṟañcāṭṭiṉēṉ}.

The accused takes -\ta{ஐ}. The -\ta{ம்} of \ta{குற்றம்} turns to \ta{ஞ்} before \ta{ச}.

\begin{cousinweb}
\textbf{\ta{குற்றஞ்சாட்டு}} (\emph{kuṟṟañcāṭṭu}) is a Tamil compound built from \ta{குற்றம்}, \textquotedblleft{}fault\textquotedblright{} (DEDR 1851), and \ta{சாட்டு}, which DEDR 2300 gives as \textquotedblleft{}to beat, strike\textquotedblright{}; how the compound came to mean to accuse is not settled.
\end{cousinweb}

\begin{grammarlens}[title={What you've built}]
One more word for telling what happened.
\end{grammarlens}

\subsection*{Your turn}
\begin{itemize}
\raggedright
  \item \emph{Say it:} \emph{kuṟṟañcāṭṭu}
  \item \emph{Say it:} \emph{kuṟṟañcāṭṭu}, once more
  \item \emph{Say it:} \emph{avar eṉṉaip poyyākak kuṟṟañcāṭṭiṉār.} — He accused me falsely.
  \item \emph{Recall:} say the Tamil for a counterfeit note, then the Tamil for an investigation, then say \emph{kuṟṟañcāṭṭu} again
\end{itemize}

\begin{tcolorbox}[breakable,colback=teal!4,colframe=teal!35!black,title={Before you move on}]
What does \ta{குற்றஞ்சாட்டு} mean? (\textquotedblleft{}To accuse\textquotedblright{}.) Say it once more.
\end{tcolorbox}
